\subsubsection{Distribución de paquetes}

\textbf{Tareas.ads}. Se definen las tareas y los recursos. Como se ha implementado concurrencia, los recursos (como la Pantalla o la Tarjeta A/D) deben protegerse para que dos tareas no escriban a la vez y mezclen las letras en la consola. Para ello se emplean \textit{entries} en tareas servidoras (Pantalla\_VP, Lector\_AD, Escritor\_AD) y un objeto protegido (Almacenamiento\_Datos) que garantiza la exclusión mutua.

\textbf{Tareas.adb}. Se implementa el comportamiento de cada tarea. Cada tarea principal tiene su propio bucle infinito. Al no haber \texttt{delay}, el procesador las alterna según la planificación que realiza el entorno de ejecución de Ada. Las tareas servidoras (\texttt{Pantalla\_VP}, \texttt{Lector\_AD}, \texttt{Escritor\_AD}) utilizan la sentencia \texttt{select} para atender \textit{entries} de forma ordenada, dando prioridad a la tarea de seguridad (\texttt{Sistema\_Seguridad}) mediante la guarda \texttt{when Mostrar\_SS'Count > 0}.

\textbf{Programa Principal concurrente.adb}. Las tareas declaradas en un paquete se activan automáticamente en cuanto el programa principal hace el \texttt{with}. Por tanto, el \texttt{main} solo tiene que quedarse esperando. Las tareas corren de forma concurrente en segundo plano, y con el \texttt{delay} se mantiene el programa principal ejecutándose.

\subsection{Resultados de la simulación}

\begin{figure}[H]
    \centering
    \includegraphics[width=0.5\linewidth]{capturas/image.png}
    \caption{Salida por consola}
    \label{fig:salida}
\end{figure}

Se muestran los mensajes de las tres tareas (CCS, CMD, SS) entrelazados en la salida. Esto indica que las tareas ocurren al mismo tiempo, compartiendo el procesador de forma concurrente.

\subsubsection{¿En qué orden se ejecutan las tareas? ¿Es siempre el mismo? Coméntalo con respecto al Ejercicio 3}

El orden de ejecución es \textbf{no determinista}. A diferencia del Ejercicio 3, donde el orden era siempre fijo (ejecución secuencial y cíclica), ahora las tareas son gestionadas por el planificador (\textit{scheduler}) del entorno de ejecución de GNAT. Las tareas se entrelazan al no haber prioridades explícitas ni retardos definidos, de modo que cada ejecución puede producir un intercalado distinto.

\subsubsection{¿Qué pasa con la tarea de seguridad? ¿En qué orden se ejecuta?}

La tarea de seguridad (\texttt{Sistema\_Seguridad}) no garantiza ser la primera en ejecutarse. Se produce una condición de carrera por ver qué tarea accede primero a los recursos compartidos. Sin embargo, en las tareas servidoras (\texttt{Pantalla\_VP} y \texttt{Lector\_AD}) se ha introducido una guarda que prioriza las peticiones de \texttt{Sistema\_Seguridad} mediante \texttt{when Mostrar\_SS'Count > 0}, lo que permite atender sus solicitudes antes que las de las otras tareas cuando hay varias pendientes. Aun así, esto no asegura que sea la primera en iniciar su bucle.

\subsubsection{¿Es necesario ahora modelar los recursos como objetos protegidos y citas? ¿Por qué?}

En este apartado es obligatorio establecer objetos protegidos y \textit{entries} para garantizar el acceso en exclusión mutua a los recursos compartidos. Al tener tres tareas concurrentes, si no usáramos estos mecanismos, dos tareas podrían intentar escribir en la pantalla al mismo tiempo, causando que las letras de una se mezclaran con las de otra en la misma línea. Las \textit{entries} en tareas servidoras permiten serializar el acceso a la pantalla y a la tarjeta A/D, mientras que el objeto protegido \texttt{Almacenamiento\_Datos} asegura que solo una tarea escriba en el almacenamiento a la vez.
